\documentclass[10pt]{article}
\usepackage[a4paper,margin=19mm]{geometry}
\usepackage{times}
\usepackage{graphicx}
\usepackage{booktabs}
\usepackage{tabularx}
\usepackage{array}
\usepackage{url}
\usepackage{hyperref}
\usepackage{amsmath}
\usepackage{float}
\usepackage{enumitem}
\setlength{\parskip}{3pt}
\setlength{\parindent}{1em}
\emergencystretch=2em
\title{Comparing Lexical and Multilingual Dense Retrieval for Hindi Document Search: An Experimental Study}
\author{Sanjit Nayyar}
\date{}
\begin{document}
\maketitle

\begin{abstract}
This study compares two methods for retrieving documents in the Hindi language; lexical retrieval (BM25) and multilingual dense retrieval (intfloat/multilingual-e5-small). The benchmark used consisted of 2210 documents, 50 Hindi queries and 244 relevance pairs. Both methods were applied to the same documents and queries for Recall@5, Recall@10 and MRR@10. BM25 achieved results of 0.1397, 0.2007 and 0.3689 whereas dense retrieval achieved results of 0.1667, 0.2167 and 0.4682. From the average of all three metrics dense retrieval performed better than lexical retrieval. The largest difference between the two was with MRR@10 at 0.0993. However, there was no statistical significance found between the two methods from the paired tests and confidence intervals with an alpha level of 0.05. Looking at the top-1 results, dense retrieval retrieved the qrels marked document 12 times and BM25 retrieved it 5 times. They both retrieved the document 8 times and neither retrieved it 25 times. This suggests that the methods are complementary to each other. The study has limitations such as the benchmark being small and the relevance judgements may be incomplete. Only one model was used for dense and lexical retrieval and no measurements were taken about how efficient they were.
\end{abstract}

\noindent\textbf{Keywords:} Hindi information retrieval; BM25; multilingual embeddings; dense retrieval; Recall; mean reciprocal rank

\section{Introduction}
Information Retrieval (IR) systems assist users in finding relevant information by interpreting a query and retrieving documents based on their relevance to the query. As users typically only view very few documents, having relevant documents ranked at the bottom of a result list may prevent users from locating them.

When trying to make an effective IR system, it can be challenging for queries and documents that are not in English. Different languages will have a different vocabularly, morphology, terminology and writing style. There will also be fewer tools available for some languages. This means that multilingual information retrieval is an important area to explore, especially for languages like Hindi.

Lexical and dense retrieval are 2 common methods. Lexical retrieval considers the terms within a query. BM25 is a popular method which uses term frequency and document length normalisation to rank the documents. Dense retrieval is when the query and documents are converted to numerical vectors and compared based on there similarity. This can be useful for understanding related meanings between queries and documents, however this will depend on the model and data the system has been trained on.

This study will compare the performances of lexical and dense retrieval. This was done using the Hindi version of Bharat-NanoBEIR. This was taken from NanoSCIDOCS. Both the queries and documents are in Hindi meaning this is a Hindi to Hindi experiment, not a Hindi to English cross lingual experiment. The chosen dataset does not include the English document collection.

\textbf{Research Question.} How does using a multilingual dense retrieval method compare to BM25 on a fixed Hindi document retrieval benchmark? Using intfloat/multilingual-e5-small, what is the difference in performance measured using Recall@5, Recall@10, MRR@10? What patterns can be seen from the top 1 results on a query level?

\textbf{Objectives.} To develop a lexical BM25 baseline and multilingual dense retrieval system. To assess both systems using the same queries, corpus and relevance judgements. To analyse the query level results to understand where the systems agree and disagree. \textbf{Contribution.} A small and repeatable comparison between the systems. An analysis of how the systems behave. No new ranking algorithm is proposed. No claims are made that one system will outperform another.

\section{Background and Related Work}
\subsection{Lexical retrieval and BM25}
BM25 is an example of a probabilistic model that combines the impact of query terms, their frequency within documents and a normalisation factor for document length. This helps prevent term saturation where repeated appearances of terms will not lead to scores increasing infinitely. Similarly it prevents long documents being favoured due to having more terms. Robertson and Zaragoza have included a detailed explanation of the probabilistic theory behind BM25 and other models . \cite{robertson2009bm25}

One advantage of using a lexical baseline is that you can understand how the model is ranking using the matching words. It is also a fairly easy method to implement and means you do not need another neural encoder. However, this can result in issues where a relevant document may not be returned if there are different ways to express the same concept. This can also cause issues where the returned document contains the same high frequency words but does not relate to the same information need. These are common issues with all term matching models and how they occur in this specific Hindi benchmark will be explored.

\subsection{Multilingual dense retrieval}
This method encodes the query and candidate passages into numerical vectors. These vectors are then ranked using a similarity function such as cosine similarity. The multilingual E5 family have been trained using multilingual text pairs and contrastive learning. The technical report provides details about these models and how they were trained/evaluated . The multilingual-e5-small model has been chosen. This is a multilingual model that outputs 384 dimensional embeddings. The model card provides information on how the model should be used. This includes asymmetric retrieval tasks using the query: and passage: prefixes. The model card also states that input longer than 512 tokens will be truncated. \cite{wang2024multilinguale5,e5modelcard}

These prefixes are part of the input format for the model. For this experiment the query prefix has been used for the search query and the passage prefix for each document in the corpus. The model card does state that performance may be worsened with lower resource languages.

\subsection{Benchmark context}
The dataset card states that the chosen resource is the Hindi version of NanoSCIDOCS inside the Bharat-NanoBEIR resource. This includes the corpus, queries and qrels linking the queries to relevant documents. The NanoBEIR datasets are small benchmarks with a small number of queries. The dataset card states that the language is Hindi and licensed under CC-BY-4.0. This makes it useful for a limited experiment however the size means it will not be a strong conclusion. \cite{bharatnanobeir2024}

\section{Dataset and Data Validation}
\subsection{Dataset composition}
The dataset was loaded into the experiment using the Hugging Face Datasets library with the id 'carlfeynman/Bharat\_NanoSCIDOCS\_hi'. The corpus, queries and qrels were loaded with the available training split. The number of documents in the corpus was 2210, there were 50 queries in total and 244 relevance judgments. Each query had between 3-5 relevant documents according to the qrels.

It also was checked that there were no discrepancies. Every query ID used in the qrels were present in the query set and every corpus ID used in the qrels were also present in the corpus. This ensures that no query or corpus IDs will be excluded due to a discrepancy which may result in loss of information we know is relevant. However, this does not necessarily mean that the relevance judgments are comprehensive.

\subsection{Interpretation of qrels}
The qrels give binary relevance information. This means they do not explicitly state that every document in the corpus is not relevant. Any document not listed in the qrels for a specific query should be stated as not being marked as relevant in the qrels. Since there are few relevance judgments given in this benchmark, it may just not have been judged. Only the positive judgments are considered when calculating the metrics.

\subsection{Preprocessing}
For BM25, every text was lower cased and tokenised using the regex w+ in Python. No stemming, lemmatisation or stopword removal was done. This is a simple way to ensure that the tokenisation can be reproduced. However it does not involve any language specific morphological analysis for Hindi. There was no further translation step. For dense retrieval, the queries and documents were prefixed with the required information for E5 and encoded using the chosen Sentence Transformers model.

\section{Methods}
\subsection{BM25 baseline}
The lexical baseline used rank\_bm25.BM25Okapi. The tokenized corpus documents were used to construct one BM25 index. For each query, the same tokenizer generated the query terms, and the BM25 implementation returned a score for each corpus document. Documents were sorted in descending score order. The notebook did not explicitly set BM25 hyperparameters, so the library defaults were used. No parameter search or relevance-feedback tuning was performed on the 50-query benchmark.

\subsection{Dense retrieval system}
The dense system used intfloat/multilingual-e5-small through the Sentence Transformers library. The notebook selected the device using torch.cuda.is\_available(); the recorded run used the CPU. Document inputs were formed as passage:  followed by the document text, and query inputs as query:  followed by the query text. Both sets of texts were encoded with batch size 32, normalized embeddings, and NumPy output. The resulting document embedding matrix had shape $2210\times384$, and the query embedding matrix had shape $50\times384$.

Because the embeddings were normalized, their dot product is equivalent to cosine similarity, subject to numerical precision. The dense score matrix was calculated as
\[
S = QD^{\mathsf T},
\]
where $Q$ is the matrix of query embeddings and $D$ is the matrix of document embeddings. For each query, the documents were sorted by descending similarity. No fine-tuning of E5 was performed on this dataset.

\subsection{Evaluation metrics}
Both methods ranked the same 2,210 documents for each of the same 50 queries. Let $R_q$ denote the set of documents marked relevant for query $q$, and let $T_{q,k}$ denote the top-$k$ ranked documents for that query. Recall at cutoff $k$ was calculated as
\[
\operatorname{Recall@k}(q)=\frac{|T_{q,k}\cap R_q|}{|R_q|}.
\]
The reciprocal rank at cutoff 10 was calculated from the rank of the first qrels-marked relevant document in the top ten:
\[
\operatorname{RR@10}(q)=\begin{cases}
1/r_q,&\text{if the first relevant result has rank }r_q\leq10,\\
0,&\text{if no relevant result occurs in the top ten.}
\end{cases}
\]
Mean reciprocal rank at 10 (MRR@10) is the arithmetic mean of these per-query reciprocal ranks. The reported Recall@5, Recall@10, and MRR@10 values are macro-averages over the 50 queries, so each query contributes equally to the mean.

\subsection{Paired query-level analysis}
For each query we have pairs of data. In addition to descriptive analysis (wins, losses and draws), an exploratory Wilcoxon signed rank test was performed to determine if there was a significant difference between the pairwise difference in each metric. Two-tailed tests were conducted. Holm correction was applied for the 3 separate metrics to correct for the family-wise error rate. A paired non-parametric bootstrap (n=50000, ran with seed = 42) was used to determine the percentile 95\% confidence intervals for the mean pairwise difference in each metric. This was a post-hoc analysis on the CSV outputs provided. It was not included in the original protocol described in the notebook.

\section{Results}
\subsection{Aggregate retrieval performance}
The macro averaged results are shown in Table~\ref{tab:main-results}. The dense system (E5-small) had a higher observed average across all 3 metrics. The greatest absolute difference occurred in the MRR@10 metric (0.0993). Recall@5 had an absolute difference of 0.0270 whilst Recall@10 had an absolute difference of 0.0160. Overall this indicates that the dense system returned the first relevant document (qrels) sooner on average. However, this does not indicate a reliable gain across a wider population of queries.

\begin{table}[H]
\centering
\caption{Macro-averaged retrieval metrics over 50 Hindi queries.}
\label{tab:main-results}
\begin{tabular}{lrrr}
\toprule
Metric & BM25 & Dense E5-small & Difference\\
\midrule
Recall@5 & 0.1397 & 0.1667 & +0.0270\\
Recall@10 & 0.2007 & 0.2167 & +0.0160\\
MRR@10 & 0.3689 & 0.4682 & +0.0993\\
\bottomrule
\end{tabular}
\end{table}

\begin{figure}[H]
\centering
\includegraphics[width=0.78\linewidth]{figures/metric_comparison.png}
\caption{Observed macro-average retrieval metrics for BM25 and dense E5-small.}
\label{fig:metric-comparison}
\end{figure}

\subsection{Query-level wins, ties, and losses}
The query-level comparison provides additional context beyond the aggregate means. For Recall@5, dense retrieval had a higher score on 12 queries, the methods tied on 30, and BM25 had a higher score on 8. For Recall@10, dense retrieval won on 12 queries, tied on 26, and BM25 won on 12. For MRR@10, dense retrieval won on 17 queries, tied on 22, and BM25 won on 11. The high number of ties, particularly for Recall@5, shows that the average difference is not a uniform improvement on every query.

\begin{table}[H]
\centering
\caption{Per-query wins, ties, and losses.}
\label{tab:win-tie-loss}
\begin{tabular}{lrrr}
\toprule
Metric & Dense wins & Ties & BM25 wins\\
\midrule
Recall@5 & 12 & 30 & 8\\
Recall@10 & 12 & 26 & 12\\
MRR@10 & 17 & 22 & 11\\
\bottomrule
\end{tabular}
\end{table}

\subsection{Top-1 qrels-membership patterns}
For each query, the first-ranked document from each method was checked for membership in that query's supplied qrels. Dense retrieval alone returned a qrels-marked relevant top-1 document for 12 queries. BM25 alone did so for 5 queries. Both systems returned a qrels-marked relevant top-1 document for 8 queries, while neither did so for 25 queries. These are qrels-membership patterns; the 25 cases in which neither top result was listed as relevant should not be interpreted as proof that both retrieved documents were irrelevant.

\begin{figure}[H]
\centering
\includegraphics[width=0.72\linewidth]{figures/top1_patterns.png}
\caption{Top-1 patterns based on membership in the supplied qrels.}
\label{fig:top1-patterns}
\end{figure}

\subsection{Exploratory paired statistical analysis}
Table~\ref{tab:paired-tests} summarizes the post-hoc paired analysis. The bootstrap intervals for the mean differences include zero for all three metrics. The two-sided Wilcoxon tests also have -values above 0.05; the Holm-adjusted values remain above 0.05. Therefore, the observed positive average differences should be treated as descriptive evidence from this benchmark, not as statistically significant improvements.

\begin{table}[H]
\centering
\caption{Paired post-hoc analysis of dense-minus-BM25 query-level differences.}
\label{tab:paired-tests}
\small
\begin{tabularx}{\linewidth}{lrrrr}
\toprule
Metric & Mean diff. & Bootstrap 95\% CI & Wilcoxon $p$ & Holm $p$\\
\midrule
Recall@5 & 0.0270 & [-0.0130, 0.0680] & 0.0832 & 0.2497\\
Recall@10 & 0.0160 & [-0.0400, 0.0720] & 0.3742 & 0.3742\\
MRR@10 & 0.0993 & [-0.0303, 0.2285] & 0.1287 & 0.2575\\
\bottomrule
\end{tabularx}
\end{table}

\begin{figure}[H]
\centering
\includegraphics[width=0.76\linewidth]{figures/paired_difference_ci.png}
\caption{Paired bootstrap percentile 95\% intervals for mean per-query differences. All intervals cross zero.}
\label{fig:paired-ci}
\end{figure}

\section{Query-Level Error Analysis}
The top-1 results have shown that the two approaches do not give identical results for all queries. A few examples have been provided in the notebook. These are not a thorough qualitative analysis of all 50 queries. They simply provide a few examples of when results differ.

For the Operational Resilience and Financial Hedging query, the top 3 results returned by BM25 were not marked as relevant in the provided qrels. However, the first result returned by the dense system was marked as relevant. This example shows that the dense system may provide a more semantically similar result. It does not prove that the dense system will always handle paraphrasing better though.

For the Home Sensors, Machine Learning and Daily Activity Patterns for Dementia query, the top 3 results returned by BM25 were not marked as relevant and the first result returned by the dense system was marked as relevant. This shows that the dense systems first result is more relevant to the main topic of the query.

However, for the Speedup Aggregation Using In-Cycle Parallelism query, BM25 returns a relevant first result about SIMD instructions. The dense systems top 3 results were not marked as relevant. This example shows how lexical retrieval can work well with specific technical terminology. The FingerCode Fingerprint query also shows that BM25 has the advantage in the top-1 flags. The Dense document snippet is however blank in the provided error analysis CSV.

Overall it seems that the dense system may provide advantages when the relevant documents contain semantically related information. Whereas lexical retrieval may benefit from specific terminology. This is based on a few examples and the experiment does not isolate what caused the documents to be ranked in this way. No human assessment of relevance was performed.

\section{Discussion}
Overall, dense E5-small achieved better averages across all three metrics however the statistical tests did not show a statistically significant improvement to the 0.05 level. This means that not all queries would benefit from having a higher average score. These results will be based on this benchmark.

Looking at results for each query shows that neither model performed well on all queries. Dense retrieval obtained more top 1 results with only relevance found through its ranking however there were instances where BM25 achieved top results with relevant documents. It appears from these examples that dense E5-small may handle technical terms and queries with semantically related concepts differently to BM25.

A combination of both models could be useful in future work. A hybrid model could use lexical retrieval to retain the exact technical terms whilst also using dense retrieval to identify similar semantics. However, no hybrid models were tested so we do not know how effective they would be.

The qrels also effect how the results are interpreted. Any document not included in the qrels is considered to not be relevant when calculating the metrics. This does not necessarily mean that they are not relevant. There may be some documents with useful information which have not been judged. Further investigation would be useful with more complete human annotations of relevance.

\section{Limitations and Threats to Validity}
\begin{enumerate}[leftmargin=*,itemsep=2pt]
\item Sample size: Only 50 queries used for evaluation. Results may vary significantly at the query level, and may not be generalisable to other Hindi collections/query distributions.
\item qrels: Only positive relevance judgements provided. No negatives provided meaning documents not identified as relevant have not been judged.
\item Not covering multiple methods: Comparison of BM25 and Dense model. Other tokenisers have not been compared, no other embedding models, no reranking or hybrid systems.
\item Basic pre-processing for Hindi: BM25 uses regex tokeniser. No stemming, lemma or stop words applied. No specific morphological techniques for Hindi tested.
\item Model truncation: Multilingual E5 model card mentions input sequences will be truncated if longer than 512 tokens. Notebook does not state how many tokens each document had or how many were truncated. Therefore cannot quantify this effect.
\item No fine-tuning/optimisation: Neither model was tuned on the benchmark. This means there was no possibility of choosing parameters based on the evaluation queries. However, this also means we are comparing the models in their chosen out-of-box form.
\item Efficiency: Notebook provides information on CPU. No efficiency comparison provided in a controlled manner. No claims of efficiency.
\item Unreliable statistics: Post-hoc paired tests and bootstrap intervals could not show a statistically significant difference. The statistics are exploratory and only 50 paired query results were used.
\item Qualitative analysis: Only specific examples investigated. Five Dense documents in the provided error analysis CSV had empty snippet fields. These have not been considered.
\end{enumerate}

\section{Reproducibility and Implementation Record}
The recorded execution environment was Google Colab on Linux x86\_64, using Python 3.13.18. The recorded package versions were rank-bm25 0.2.2, datasets 4.8.5, transformers 5.22.1, tokenizers 0.23.3, sentence-transformers 5.7.0, torch 2.14.1+cpu, numpy 2.4.6, and pandas 2.2.3. The run used the CPU, and the execution timestamp printed by the notebook was 2026-11-13 08:45:49 UTC.

These package versions describe the recorded run; they do not guarantee that a future Colab session will recreate the same environment because the installation commands were not pinned. The ZIP supplied for this paper contains the LaTeX source, bibliography, and project documentation. It does not contain the executed notebook or the query-level CSV result files. Therefore, this archive alone is not sufficient to rerun the experiment. The statistical analysis reported in this paper was computed separately from the supplied query-level CSV outputs using paired Wilcoxon signed-rank tests and a percentile bootstrap with 50,000 resamples and random seed 42.

\section{Conclusion and Future Work}
For this study, BM25 was compared with multilingual dense retrieval (intfloat/multilingual-e5-small). Dense retrieval outperformed BM25 on average in Recall@5, Recall@10 and MRR@10 scores for this static benchmark. Dense retrieval outperformed BM25 most in MRR@10; however, confidence intervals did include zero for each of these variables. This means that the results only indicate that there is a difference between the averages for the variables, but do not indicate that dense retrieval is a reliable improvement.

Looking at the results for each query it can be seen that both methods performed well but also poorly. Dense retrieval alone had 12 queries where the top 1 document was relevant (qrels marked), while BM25 alone had 5. 8 queries were deemed relevant by both methods, and 25 queries were not deemed relevant by either method. These results are based on the presence within qrels and should not be taken to mean that the documents not listed are not relevant.

In future, a larger and more diverse benchmark could be used for Hindi. More human relevance marking could be added, and other lexical and dense models could be trialled. It may also be useful to look into Hindi specific text processing. Testing how much time and memory each method takes to run in a controlled environment could also be useful, along with testing a combined approach to lexical and dense retrieval. This will help to determine if the trends seen in this experiment can be seen across other datasets.

\begin{thebibliography}{9}
\bibitem{robertson2009bm25}
S. E. Robertson and H. Zaragoza.
\newblock The probabilistic relevance framework: BM25 and beyond.
\newblock \emph{Foundations and Trends in Information Retrieval}, 3(4):333--389, 2009.
\newblock doi: \url{https://doi.org/10.1561/1500000019}.

\bibitem{wang2024multilinguale5}
L. Wang, N. Yang, X. Huang, L. Yang, R. Majumder, and F. Wei.
\newblock Multilingual E5 text embeddings: A technical report.
\newblock arXiv:2402.05672, 2024. \url{https://arxiv.org/abs/2402.05672}.

\bibitem{bharatnanobeir2024}
Bharat-NanoBEIR.
\newblock Bharat-NanoBEIR: Indian language information retrieval datasets.
\newblock Hugging Face Datasets, 2024. \url{https://huggingface.co/datasets/carlfeynman/Bharat_NanoSCIDOCS_hi}.

\bibitem{e5modelcard}
intfloat.
\newblock multilingual-e5-small model card.
\newblock Hugging Face, 2024. \url{https://huggingface.co/intfloat/multilingual-e5-small}.
\end{thebibliography}
\end{document}
